\section{Self-Distilled Admission}
\label{sec:method}

SDA distills a frozen model's candidate rankings into its answer to the
importance prompt (Figure~\ref{fig:sda}; Algorithm~\ref{alg:sda}).

\begin{figure*}[t]
\centering
\figtodo{6.2cm}
\caption{SDA. (a) The frozen model combines workspace
readout and output likelihood to rank each episode's candidates and
labels the top two \emph{yes}. (b) A LoRA adapter learns these labels
through the model's own importance prompt. (c) The serving engine applies
the adapter only to admission requests.
\textit{Drawing guide (6.3\,in $\times$ 2.45\,in, Figure~\ref{fig:hero} palette): three panels.
(a) Teacher: passage $x$ and candidates $C$ enter a frozen model (grey box with a snowflake). Beside it,
a small heat strip with layers on the vertical axis ($1$ at the bottom, $L$ at the top) and one column per
candidate, shaded by the reciprocal rank of the candidate's first token at that layer; a navy dot marks
each column's peak, labeled $W(c;x)$. The output head is labeled $\mathrm{LL}(c;x)$ in teal. Both
become within-episode percentile bars, averaged into $T$; a sorted candidate list tags the top two
\emph{yes} (rust) and the rest \emph{no} (grey). (b) Student: the same frozen model with a small rust
LoRA block, answering the importance prompt ``Is $c$ one of the most important things to remember\ldots?
yes/no''; a dashed line to the labels of (a) reads ``next-token cross-entropy''. (c) Deployment: a host
agent sends passage and candidates to a serving engine (vLLM); admission requests run with the adapter
on and store the top $k$ in memory; question requests run with the adapter off, and the frozen reader
answers.}}
\label{fig:sda}
\end{figure*}

\subsection{Teacher Rankings}
\label{sec:teacher}

The teacher ranks each episode's candidates using two signals from the
frozen model, defined next.

\paragraph{Workspace readout.}
The workspace readout measures how readily the unembedding recovers a
concept during context processing. SDA uses only the logit lens, which
decodes each layer's residual state through final normalization and
unembedding and is the J-lens with $J=I$
\citep{nostalgebraist2020logitlens,anthropic2026workspace}.
Let $h_\ell(x)\in\mathbb{R}^d$ be the residual state at the final context
position after transformer layer $\ell$, with $L$ layers in total.
Final normalization $\mathrm{LN}$ and unembedding $W_U$ give vocabulary
logits $z_\ell(x)=W_U\,\mathrm{LN}(h_\ell(x))$.
Writing $\mathrm{rank}_\ell(t)=1+|\{v:z_{\ell,v}>z_{\ell,t}\}|$, we take
the peak reciprocal rank across layers and first-token forms:
\begin{equation}
W(c;x)=\max_{\substack{1\leq\ell\leq L\\ t\in T(c)}}
\frac{1}{\mathrm{rank}_\ell(t)}.
\label{eq:wrr}
\end{equation}
Here $T(c)$ contains the distinct first tokens of the candidate's bare
and leading-space forms, each in original and capitalized form.
The maximum captures the candidate's strongest decoding rank across
these layers and forms. Computing the readout requires one forward pass
over $x$, using final-position states from all layers.

\paragraph{Combined teacher.}
Output likelihood scores the complete candidate by its mean log
probability per token:
$\mathrm{LL}(c;x)=\log p_{\mathcal M}(c\mid x)/|c|$, where $|c|$ is the
candidate's token length. To put workspace readout and output likelihood
on a common scale, we map each to mid-rank percentiles within the
episode. Tied values receive their mean rank
(App.~\ref{app:training}). The teacher used in the main evaluation
averages these percentiles:
\begin{equation}
T(c;x)=\tfrac12\bigl(P_W(c;x)+P_{\mathrm{LL}}(c;x)\bigr).
\label{eq:teacher}
\end{equation}
Here $P_S$ denotes the percentile under score $S$.
The two highest-ranked candidates receive the binary label $\tilde y=1$
for \emph{yes}; all others receive $\tilde y=0$ for \emph{no}.
Candidate order breaks score ties. The frozen model generates every
teacher label from the passage and candidates.

\subsection{Admission Distillation}
\label{sec:student}

The student is a LoRA adapter on the frozen model that learns the
teacher's labels by answering the importance prompt.

\paragraph{Admission score.}
The importance prompt $u(c,x)$ asks whether $c$ is one of the most
important concepts to remember for possible future questions and
requests a single yes or no (App.~\ref{app:prompts}). Qwen3 uses
thinking-disabled prompts. Let $Y,N$ contain the distinct first tokens
of the case and leading-space forms of \emph{yes} and \emph{no},
respectively. The admission score compares their total probabilities:
\begin{equation}
\begin{split}
s_\theta(c;x)={}&\log\!\sum_{t\in Y}p_\theta(t\mid u(c,x))\\
&-\log\!\sum_{t\in N}p_\theta(t\mid u(c,x)),\\
V_\theta(c;x)={}&\sigma\!\left(s_\theta(c;x)\right).
\end{split}
\label{eq:score}
\end{equation}
The untrained model's score $V$ is the importance report; SDA ranks
candidates by the adapter-enabled score $s_\theta$.

\paragraph{Training objective.}
We optimize only the LoRA adapter; base parameters remain frozen.
Cross-entropy trains next-token prediction of the teacher's target word:
\begin{equation}
\mathcal L(\theta)=-\frac{1}{|\mathcal D|}
\sum_{(x,c)\in\mathcal D}
\log p_\theta\!\left(\tau(\tilde y(c;x))\mid u(c,x)\right).
\label{eq:loss}
\end{equation}
The training examples in $\mathcal D$ pair each passage $x$ with each of
its candidates $c$ and the corresponding teacher label $\tilde y(c;x)$.
Here $\tau$ maps a positive label to \emph{yes} and a negative label to
\emph{no}. The target probability is normalized over the vocabulary.
Prompt tokens are masked from the loss
(App.~\ref{app:configuration}).

\paragraph{Training and selection.}
We train a rank-16 adapter for two epochs and select the earliest
checkpoint with the highest within-episode AUC on the in-distribution
validation split. Optimization uses teacher labels; validation uses
gold utility labels. We select the combined teacher by mean answer
accuracy across development sets and training seeds, scored by the
development run's automatic grader. The configuration is fixed before
evaluation on the confirmation set (App.~\ref{app:freeze}).

\begin{algorithm}[t]
\caption{Self-Distilled Admission.}
\label{alg:sda}
\small
\begin{algorithmic}[1]
\Require frozen model $\mathcal M$; training episodes $\mathcal E$, each a passage $x$ with candidates $C$
\For{each episode $(x,C)\in\mathcal E$}
  \For{each candidate $c\in C$}
    \State $W(c;x)\gets$ peak reciprocal rank across layers \Comment{Eq.~\eqref{eq:wrr}}
    \State $\mathrm{LL}(c;x)\gets\log p_{\mathcal M}(c\mid x)/|c|$
  \EndFor
  \State $T(c;x)\gets\tfrac12\bigl(P_W(c;x)+P_{\mathrm{LL}}(c;x)\bigr)$ for all $c\in C$ \Comment{Eq.~\eqref{eq:teacher}}
  \State $\tilde y(c;x)\gets 1$ for the two highest $T$, else $0$
\EndFor
\State train adapter $\theta$ on $\mathcal L(\theta)$ over all $(x,c,\tilde y)$ \Comment{Eq.~\eqref{eq:loss}}
\State select the earliest checkpoint with the highest validation within-episode AUC
\Statex \textbf{Deployment:} given a passage $x$, candidates $C$ and budget $k$
\State store the $k$ candidates with the highest $s_\theta(c;x)$ \Comment{Eq.~\eqref{eq:score}}
\end{algorithmic}
\end{algorithm}

\subsection{Off-the-Shelf Serving}
\label{sec:serving}

At deployment, SDA makes one yes/no forward pass per candidate and reads
the importance prompt's output probabilities to compute $s_\theta$.
The host stores the top-$k$ candidates and manages retrieval.
Candidate order breaks score ties.

The adapter is attached only to admission requests in an off-the-shelf
serving engine such as vLLM \citep{kwon2023pagedattention}.
Prefix caching reuses the shared passage. All other requests,
including the frozen reader's answer generation, use the unmodified
model.
